\chapter*{Zusammenfassung}
\addcontentsline{toc}{chapter}{Zusammenfassung}

\selectlanguage{german}

Graph Neural Networks werden zunehmend für klinische Vorhersagen aus
elektronischen Patientenakten eingesetzt, ihre Entscheidungen werden jedoch
meist erst im Nachhinein erklärt, durch Post-hoc-Attributionsverfahren, deren
Ausgabe sich nicht daran prüfen lässt, wie das Modell tatsächlich entschieden
hat. Diese Arbeit untersucht, ob eine selbsterklärende Architektur in diesem
Umfeld eine brauchbare Alternative darstellt. Aus MIMIC-IV-ED wird pro Besuch in
der Notaufnahme ein heterogener Graph gebildet, mit Vitalparametern,
Medikamenten, Symptomen und Aufnahmegrund als typisierten Knoten. Vorhergesagt
wird die Hauptdiagnose über 30 zusammengefasste Krankheitsklassen. Die
Vorhersage entsteht in einer Prototyp-Schicht, die jeden Patientengraphen mit 90
gelernten Archetypen vergleicht; diese werden auf reale Subgraphen aus den
Trainingsdaten projiziert. Das Modell wird mit einem Knowledge-Graph-GNN
(GraphCare), einer Ablation desselben Encoders ohne Prototyp-Schicht und einem
tabellarischen Gradient-Boosting-Verfahren verglichen, alle auf demselben
subjektbewussten Split.

Drei Ergebnisse stechen hervor. Das tabellarische Verfahren sagt in allen
Metriken am besten vorher (65,4\,\% Accuracy, 0,607 Macro-F1) gegenüber
56,3\,\% und 0,508 beim Prototyp-Modell. Die hier verwendete Graphstruktur
bringt also Interpretierbarkeit statt Genauigkeit. Die Prototyp-Schicht selbst
ist nahezu kostenlos: Wird sie entfernt, ändert sich die Accuracy um 0,2
Prozentpunkte. Und \caveat{bei einer Stichprobe von 50 Graphen des Test-Folds} wählen die drei
Post-hoc-Verfahren nur bei 32\,\% der Graphen denselben wichtigsten Knoten,
während die Prototyp-Evidenz in 39 dieser 50 Fälle zur Vorhersage passt und ihr
Abstand richtige von falschen Vorhersagen trennt. Das spricht dafür, die Erklärung im Modell selbst zu
verankern.

\tbd{Ein drittes Vergleichsverfahren (Methode~C) ist vorgesehen, aber noch nicht
implementiert. Sobald es vorliegt, hier einen Satz zu seinem Beitrag und seinen
Ergebnissen ergänzen und die Zahlen oben anpassen, falls sich die Reihenfolge
ändert.}

\selectlanguage{english}
